\section{Landasan Konsep Finansial}

% Slide 4.1: Bab 2.1.1 - Evolusi Uang
\begin{frame}{Bab 2.1.1 - Evolusi Uang}
    \begin{itemize}
        \item Transformasi dari kebutuhan biologis (Barter) $\rightarrow$ Uang Komoditas (Logam) $\rightarrow$ Uang \textit{Fiat} $\rightarrow$ Entitas Data Digital.
    \end{itemize}
    \begin{figure}
        \centering
        \begin{subfigure}[b]{0.22\textwidth}
            \centering
            
\includegraphics[width=\textwidth]{211_Barter_marketeers.com.png}
            \caption{\textit{Barter}}
        \end{subfigure}
        \hfill
        \begin{subfigure}[b]{0.22\textwidth}
            \centering
            
\includegraphics[width=\textwidth]{211_Logam_depositphotos.com.png}
            \caption{\textit{Commodity}}
        \end{subfigure}
        \hfill
        \begin{subfigure}[b]{0.22\textwidth}
            \centering
            
\includegraphics[width=\textwidth]{211_Fiat_depositphotos.com.png}
            \caption{\textit{Fiat}}
        \end{subfigure}
        \hfill
        \begin{subfigure}[b]{0.22\textwidth}
            \centering
            \includegraphics[width=\textwidth]{211_Digital_klasika.kompas.id.png}
            \caption{\textit{Digital}}
        \end{subfigure}
        \caption{Tahapan evolusi instrumen pertukaran nilai dalam peradaban manusia.}
    \end{figure}
\end{frame}

% Slide 4.2: Bab 2.1.2 - Inflasi
\begin{frame}{Bab 2.1.2 - Inflasi}
    \begin{columns}
        \begin{column}{0.5\textwidth}
            \begin{itemize}
                \item Fenomena sistemik peningkatan entropi ekonomi.
                \item Degradasi daya beli yang memaksa realokasi modal strategis untuk mempertahankan kekayaan riil.
            \end{itemize}
        \end{column}
        \begin{column}{0.5\textwidth}
            \begin{figure}
                \centering
                \includegraphics[width=\textwidth]{212_Inflasi.png}
                \caption{Visualisasi dampak inflasi terhadap daya beli.}
            \end{figure}
        \end{column}
    \end{columns}
\end{frame}

% Slide 4.3: Bab 2.1.3 - Investasi
\begin{frame}{Bab 2.1.3 - Investasi}
    \begin{columns}
        \begin{column}{0.5\textwidth}
            \begin{itemize}
                \item Pengorbanan konsumsi jangka pendek untuk imbal hasil masa depan.
                \item Prinsip \textit{Risk-Return Trade-off}: Imbal hasil tinggi selalu diiringi risiko yang setara.
            \end{itemize}
        \end{column}
        \begin{column}{0.5\textwidth}
            \begin{figure}
                \centering
                \includegraphics[width=0.8\textwidth]{213_Investasi.png}
                \caption{Spektrum risiko dan imbal hasil aset investasi.}
            \end{figure}
        \end{column}
    \end{columns}
\end{frame}

% Slide 4.4: Bab 2.1.4 - Penghindaran Risiko (Risk Aversion)
\begin{frame}{Bab 2.1.4 - Penghindaran Risiko (Risk Aversion)}
    \begin{columns}
        \begin{column}{0.5\textwidth}
            \begin{itemize}
                \item Parameter endogen yang dipengaruhi oleh Teori Prospek.
                \item Respon asimetris terhadap kerugian (\textit{loss aversion}) dibandingkan keuntungan.
            \end{itemize}
        \end{column}
        \begin{column}{0.5\textwidth}
            \begin{figure}
                \centering
                \includegraphics[width=\textwidth]{214_RiskAversion_invespcro.com.png}
                \caption{Fungsi Utilitas dalam Teori Prospek.}
            \end{figure}
        \end{column}
    \end{columns}
\end{frame}

% Slide 5.1: Bab 2.2.1 - Analisis Kualitatif (Fundamental)
\begin{frame}{Bab 2.2.1 - Analisis Kualitatif (Fundamental)}
    \begin{columns}
        \begin{column}{0.5\textwidth}
            \begin{itemize}
                \item Evaluasi kesehatan industri dan emiten secara menyeluruh.
                \item Fokus pada laporan keuangan (\textit{Cashflow statement}), neraca, dan sentimen pasar.
                \item Menilai nilai intrinsik aset berdasarkan data makro dan mikro ekonomi.
            \end{itemize}
        \end{column}
        \begin{column}{0.5\textwidth}
            \begin{figure}
                \centering
                \includegraphics[width=\textwidth]{221_Cashflow_statement.png}
                \caption{Komponen krusial dalam analisis fundamental.}
            \end{figure}
        \end{column}
    \end{columns}
\end{frame}

% Slide 5.2: Bab 2.2.2 - Analisis Teknikal (Visual)
\begin{frame}{Bab 2.2.2 - Analisis Teknikal (Visual)}
    \begin{itemize}
        \item Identifikasi pola memori kolektif pelaku pasar melalui representasi visual harga.
    \end{itemize}
    \begin{figure}
        \centering
        \begin{subfigure}[b]{0.3\textwidth}
            \centering
            \includegraphics[width=\textwidth]{222_Candle.png}
            \caption{\textit{Candlestick}}
        \end{subfigure}
        \hfill
        \begin{subfigure}[b]{0.3\textwidth}
            \centering
            \includegraphics[width=\textwidth]{222_Chart.png}
            \caption{\textit{Price Chart}}
        \end{subfigure}
        \hfill
        \begin{subfigure}[b]{0.3\textwidth}
            \centering
            \includegraphics[width=\textwidth]{222_SnR.png}
            \caption{\textit{Support \& Resistance}}
        \end{subfigure}
        \caption{Instrumen utama dalam pemetaan psikologi pasar secara teknikal.}
    \end{figure}
\end{frame}

% Slide 5.3: Bab 2.2.3 - Analisis Kuantitatif
\begin{frame}{Bab 2.2.3 - Analisis Kuantitatif}
    \begin{columns}
        \begin{column}{0.5\textwidth}
            \begin{itemize}
                \item Transformasi harga mentah menjadi \textit{Log-Return} untuk mencapai stasioneritas data.
                \item Penting untuk pemodelan statistik dan input algoritma optimasi.
                \item Normalisasi rentang nilai guna kompatibilitas dengan sirkuit kuantum.
            \end{itemize}
        \end{column}
        \begin{column}{0.5\textwidth}
            \begin{figure}
                \centering
                \includegraphics[width=\textwidth]{223_Price_log.png}
                \caption{Transformasi pergerakan harga menjadi deret imbal hasil (\textit{returns}).}
            \end{figure}
        \end{column}
    \end{columns}
\end{frame}

% Slide 5.4: Bab 2.2.4 - Optimasi Portofolio Markowitz
\begin{frame}{Bab 2.2.4 - Optimasi Portofolio Markowitz}
    \begin{columns}
        \begin{column}{0.5\textwidth}
            \begin{itemize}
                \item \textbf{Mean-Variance Optimization:} Mencari bobot aset ($\omega$) yang meminimalkan risiko untuk target imbal hasil tertentu.
                \item \textbf{Efficient Frontier:} Himpunan portofolio optimal yang menawarkan imbal hasil maksimal pada tingkat risiko tertentu.
            \end{itemize}
        \end{column}
        \begin{column}{0.5\textwidth}
            \begin{figure}
                \centering
                \includegraphics[width=\textwidth]{224_Efficient_Frontier.png}
                \caption{Kurva batas efisiensi alokasi aset.}
            \end{figure}
        \end{column}
    \end{columns}
\end{frame}
